\chapter{Buông và giữ}
\markboth{Buông và giữ}{Release and Keep}

\section{Một người giữ mãi cơn giận nên mất cả sự nhẹ lòng.}
\begin{truyen}{Một người giữ mãi cơn giận nên mất cả sự nhẹ lòng.}{Holding hurts the holder too.}
\chuhoa{C}{àng nghĩ lại điều cũ, anh càng thấy mình đúng. Nhưng cùng lúc, anh cũng thấy mình ngày càng nặng nề. Có những thứ giữ lại không làm ta mạnh hơn, chỉ làm ta mệt hơn.}
\textit{(The more he replayed the old hurt, the more right he felt. Yet he also became heavier inside. Some things we hold do not make us stronger, only more tired.)}
\end{truyen}
\begin{baihoc}
Buông không phải luôn là thua; nhiều khi là tự cứu mình khỏi cái nặng không đáng giữ.
\textit{(Letting go is not always losing; often it is saving yourself from unnecessary weight.)}
\end{baihoc}
\begin{ghinhoanh}\textbf{Short English to remember:}

Holding hurts the holder too.
\end{ghinhoanh}
\begin{tuhocnhanh}
1. Điều gì em đang giữ mà làm em mệt? \textit{(What are you holding that is tiring you?)}

2. Em có thể buông bớt điều gì? \textit{(What can you loosen your grip on?)}
\end{tuhocnhanh}
\ngancach

\section{Một người buông tự trọng để giữ một mối quan hệ.}
\begin{truyen}{Một người buông tự trọng để giữ một mối quan hệ.}{Not everything worth keeping is worth keeping at all costs.}
\chuhoa{V}{ì sợ mất người kia, cô chấp nhận bị xem thường hết lần này tới lần khác. Đến khi còn lại rất ít chính mình, cô mới hiểu có những thứ giữ bằng cách mất bản thân thì không còn là giữ nữa.}
\textit{(Afraid of losing someone, she accepted disrespect again and again. When very little of herself remained, she understood that if keeping something costs your selfhood, it is no longer truly keeping.)}
\end{truyen}
\begin{baihoc}
Giữ người mà mất mình thường là một kiểu mất lớn hơn vẻ ngoài của nó.
\textit{(Keeping someone while losing yourself is often a greater loss than it appears.)}
\end{baihoc}
\begin{ghinhoanh}\textbf{Short English to remember:}

Do not keep others by losing yourself.
\end{ghinhoanh}
\begin{tuhocnhanh}
1. Em có đang đánh đổi lòng tự trọng để giữ gì đó không? \textit{(Are you trading self-respect to keep something?)}

2. Em muốn giữ gì mà vẫn còn là mình? \textit{(What do you want to keep without losing yourself?)}
\end{tuhocnhanh}
\ngancach

\section{Một người buông thói quen xấu rồi mới thấy mình được lại tự do.}
\begin{truyen}{Một người buông thói quen xấu rồi mới thấy mình được lại tự do.}{Giving up can also be gaining back.}
\chuhoa{B}{ỏ một thói quen tưởng như mất vui, nhưng vài tháng sau anh thấy đầu óc sáng hơn, sức khỏe khá hơn, và lòng không còn bị kéo đi mãi. Lúc đó anh hiểu buông bỏ cũng là một cách được lại.}
\textit{(Quitting a bad habit first felt like losing pleasure, but months later his mind was clearer, health better, and heart less dragged around. Then he understood that letting go can also be getting back.)}
\end{truyen}
\begin{baihoc}
Không phải thứ gì rời khỏi đời mình cũng để lại khoảng trống; có thứ rời đi để trả lại khoảng sáng.
\textit{(Not everything that leaves your life creates emptiness; some things leave to return light.)}
\end{baihoc}
\begin{ghinhoanh}\textbf{Short English to remember:}

Letting go can be gaining back.
\end{ghinhoanh}
\begin{tuhocnhanh}
1. Điều gì em bỏ đi thì đời sẽ nhẹ hơn? \textit{(What could you let go of to make life lighter?)}

2. Em đang sợ mất cảm giác gì? \textit{(What feeling are you afraid to lose?)}
\end{tuhocnhanh}
\ngancach

\section{Một người giữ lời hứa nhỏ nên mất một tiện lợi trước mắt.}
\begin{truyen}{Một người giữ lời hứa nhỏ nên mất một tiện lợi trước mắt.}{Integrity often looks inconvenient first.}
\chuhoa{A}{nh có thể bẻ lời một chút để đỡ phiền, nhưng cuối cùng vẫn chọn giữ đúng điều đã nói. Cái được ngay trước mắt mất đi, nhưng lòng tự trọng ở lại.}
\textit{(He could have bent his word a little for convenience, but chose to keep it. The immediate gain disappeared, but self-respect remained.)}
\end{truyen}
\begin{baihoc}
Điều ngay tiện chưa chắc quý hơn điều lâu bền trong nhân cách.
\textit{(What is immediately convenient is not always more valuable than long-term character.)}
\end{baihoc}
\begin{ghinhoanh}\textbf{Short English to remember:}

Integrity is worth small inconvenience.
\end{ghinhoanh}
\begin{tuhocnhanh}
1. Em có đang định dễ dãi với lời hứa nào không? \textit{(Are you about to be careless with any promise?)}

2. Em muốn mình là người thế nào khi không ai ép? \textit{(What kind of person do you want to be when no one forces you?)}
\end{tuhocnhanh}
\ngancach

\section{Một người buông sĩ diện nên được lại bạn bè.}
\begin{truyen}{Một người buông sĩ diện nên được lại bạn bè.}{Pride is expensive; humility repairs.}
\chuhoa{S}{au thời gian dài im lặng, anh chủ động mở lời trước. Chỉ một bước xuống khỏi sĩ diện đã kéo lại một mối quan hệ tưởng mất rồi.}
\textit{(After a long silence, he chose to speak first. One step down from pride brought back a relationship that seemed lost.)}
\end{truyen}
\begin{baihoc}
Có khi thứ ta tưởng mất không phải vì đã muộn, mà vì chưa ai chịu hạ cái tôi xuống đủ thấp.
\textit{(Sometimes what we think is lost is not gone because it is too late, but because no one has lowered their ego enough.)}
\end{baihoc}
\begin{ghinhoanh}\textbf{Short English to remember:}

Humility can recover what pride lost.
\end{ghinhoanh}
\begin{tuhocnhanh}
1. Sĩ diện đang chặn em quay lại với ai? \textit{(How is pride stopping you from returning to someone?)}

2. Một lời mở đầu có thể là gì? \textit{(What could your opening sentence be?)}
\end{tuhocnhanh}
\ngancach